\subsection{Objective}

In Phase 2, we use the phase difference between two antennas to estimate the
direction of arrival (DoA) of a signal. The current program uses two receive
channels to estimate this angle (\texttt{decode.py}). It is a first step
towards beamforming and source tracking. It is not yet a complete system.

\subsection{Principle of phase-based direction estimation}

Consider two antennas with a distance \(d\) between them. We measure the
arrival angle \(\theta\) from broadside. Broadside is the direction
perpendicular to the line between the antennas, so \(\theta=0\) means the
wave arrives from that direction. The difference in the distances travelled
by the wave is
\[
    \Delta \ell = d\sin(\theta).
\]
For a signal with wavelength \(\lambda\), this distance difference creates a
phase difference
\[
    \Delta\phi
    = \frac{2\pi}{\lambda}\Delta\ell
    = \frac{2\pi d}{\lambda}\sin(\theta),
    \qquad
    \lambda = \frac{c}{f},
\]
Here, \(f\) is the signal frequency and \(c\) is the speed of light. We can
use the measured phase difference to estimate the angle:
\[
    \theta
    = \arcsin\!\left(\frac{\Delta\phi\,\lambda}{2\pi d}\right).
\]
The measured phase repeats every \(2\pi\). Because of this, the estimate also
depends on the antenna spacing, calibration and angle range.

\subsection{Two-channel acquisition and signal processing}

The program sets up two receive channels and reads both signals at the same
time. It uses a \(2.437~\mathrm{GHz}\) signal, a sample rate of
\(2~\mathrm{MSample/s}\) and buffers of \(16{,}384\) samples
(\texttt{decode.py}). The receiver is tuned \(200~\mathrm{kHz}\) below the
expected signal frequency. This moves the tone away from zero frequency in
the sampled signal.

For each buffer, the program applies a Hann window and computes an FFT for
each channel. It then finds the strongest frequency bin in the selected band.
Let \(X_1[k]\) and \(X_2[k]\) be the complex FFT values from channels 1 and 2
at bin \(k\). We multiply one value by the complex conjugate of the other:
\[
    C[k] = X_2[k]X_1^{*}[k],
\]
Here, \(X_1^{*}[k]\) is the complex conjugate of \(X_1[k]\). The phase
difference is the angle of this product:
\[
    \Delta\phi[k]
    = \arg\!\left(C[k]\right)
    = \arg\!\left(X_2[k]\right)-\arg\!\left(X_1[k]\right).
\]
The program repeats this calculation for several buffers and averages the
results. For \(M\) buffers, the average is
\[
    \overline{C}
    = \frac{1}{M}\sum_{m=1}^{M}
      X_{2,m}[k_m]X_{1,m}^{*}[k_m],
\]
The phase estimate is the angle of this average. Averaging helps reduce noise.

At startup, the program averages 50 measurements to get a calibration value.
During operation, it averages 10 measurements at a time, removes the
calibration offset and smooths the phase estimate. It then uses the wavelength
and antenna spacing to calculate the angle. If the result is outside the
valid range for the inverse sine, the program reports that the angle is out
of range.

\subsection{Effect of antenna spacing}

At \(2.437~\mathrm{GHz}\), the wavelength is about \(0.123~\mathrm{m}\). The
program uses an antenna spacing of \(0.0615~\mathrm{m}\), about half a
wavelength (\texttt{decode.py}). A larger spacing gives a larger phase
change for the same change in angle, which can make the measurement more
sensitive. But if the spacing is too large, different angles can give the
same measured phase. This makes the angle unclear. A spacing no greater than
half a wavelength helps avoid this problem over the main angle range.

We also need to know the spacing accurately and calibrate both receive
channels. Cables, antennas, receiver differences and reflections can all
change the measured phase. The program corrects a fixed phase difference
between the channels. We still need to choose and test a known reference
direction.

\subsection{Code example}



Le firmware de l'ESP32, il est où, wesh ?

\subsection{Results}